\documentclass[11pt,a4paper]{article}
\usepackage[utf8]{inputenc}
\usepackage[T1]{fontenc}
\usepackage[italian]{babel}
\usepackage[a4paper,top=1.8cm,bottom=1.8cm,left=2cm,right=2cm,headheight=15pt]{geometry}
\usepackage{amsmath,amssymb}
\usepackage{graphicx}
\usepackage{tikz}
\usepackage{pgffor}
\usepackage{enumitem}
\usepackage{fancyhdr}
\usepackage{lastpage}
\usepackage{booktabs}
\usepackage{longtable}
\usepackage{siunitx}
\usepackage{xcolor}
\usepackage{needspace}
\usepackage{array}
\usepackage{textcomp}
\usepackage{pdflscape}
\usepackage[hidelinks]{hyperref}
\sisetup{output-decimal-marker={,},group-separator={.},per-mode=symbol}
\setlength{\parindent}{0pt}
\setlength{\parskip}{2pt}
\setlist[enumerate]{itemsep=6pt,topsep=4pt,leftmargin=8mm}
\setlist[itemize]{itemsep=2pt,topsep=3pt,leftmargin=6mm}
\raggedbottom

% Comandi personalizzati: modifica qui la grafica senza toccare il codice.
\pagestyle{fancy}
\fancyhf{}
\renewcommand{\headrulewidth}{0.4pt}
\fancyfoot[C]{\small Pag. \thepage\ di \pageref{LastPage}}

% Punteggio a destra della consegna.
\newcommand{\punti}[1]{\hfill\textbf{[#1~pt]}}

% Riga puntinata per lo svolgimento.
\newcommand{\riga}{\noindent\makebox[\linewidth]{\dotfill}\par\vspace{4mm}}

% N righe puntinate (se N = 0 non stampa nulla: il chiamante evita il caso).
\newcommand{\rispostarighe}[1]{\par\vspace{1mm}\foreach \i in {1,...,#1}{\riga}}

% Riquadro per la risposta.
\newcommand{\riquadro}{\par\noindent\framebox[\linewidth]{\rule{0pt}{2.6cm}}\par}

% Griglia quadrettata per grafici e disegni.
\newcommand{\griglia}{\par\noindent
  \begin{tikzpicture}[x=0.55cm,y=0.55cm]
    \draw[step=1,lightgray,very thin] (0,0) grid (16,10);
    \draw[black!60,thick] (0,0) rectangle (16,10);
  \end{tikzpicture}\par}

% Piano cartesiano con assi e scala.
\newcommand{\grigliacartesiana}{\par\noindent
  \begin{tikzpicture}[x=0.55cm,y=0.55cm]
    \draw[step=1,lightgray,very thin] (0,0) grid (16,10);
    \draw[->,thick] (0,5) -- (16,5) node[right] {\small $x$};
    \draw[->,thick] (8,0) -- (8,10) node[above] {\small $y$};
    \foreach \x in {2,4,6,10,12,14} \draw (\x,4.8) -- (\x,5.2);
    \foreach \y in {1,3,7,9} \draw (7.8,\y) -- (8.2,\y);
  \end{tikzpicture}\par}

% Intestazione della verifica: istituto, indirizzo, materia, classe, data, tempo, punti, soglia.
\newcommand{\intestazione}[8]{%
  \begin{center}
    {\large\bfseries #1}\\[1mm]
    {\small #2\ --\ #3}\\[3mm]
    {\Large\bfseries #4}\\[3mm]
    \begin{tabular}{@{}p{0.48\textwidth}p{0.48\textwidth}@{}}
      Classe: #5 & Data: #6 \\
      Tempo: #7 & Punteggio: #8 \\
    \end{tabular}
  \end{center}
}

% Titolo di sezione con spazio garantito.
\newcommand{\titolosezione}[1]{\par\needspace{4\baselineskip}\section*{#1}\vspace{-1mm}}

% Blocco soluzione.
\newcommand{\soluzione}[1]{\par\noindent #1\par}

% Casella per gli errori comuni (soluzioni).
\newcommand{\erroricomuni}[1]{%
  \par\smallskip
  \noindent\fcolorbox{black!25}{black!3}{%
    \begin{minipage}{\dimexpr\linewidth-2\fboxsep-2\fboxrule\relax}
      \small\textbf{Errori comuni attesi:} #1
    \end{minipage}}\par\smallskip
}

\fancyhead[L]{\small Liceo Statale di Esempio}
\fancyhead[R]{\small Soluzioni\ --\ Matematica}

\begin{document}

\begin{center}
  {\Large\bfseries Soluzioni}\\[2mm]
  {\large Verifica dimostrativa --- Equazioni e geometria (1\textsuperscript{a} Classico)}\\[1.5mm]
  {\small Liceo Classico\ --\ Classe 1\textsuperscript{a} A\ --\ 12/11/2026\ --\ 100 punti}
\end{center}
\vspace{1mm}
\noindent\rule{\linewidth}{0.4pt}



\titolosezione{Quesiti a risposta chiusa}

\needspace{9\baselineskip}
\subsection*{Quesito 1\hfill\normalsize\mdseries [4,5~pt]}

{\small \textit{Argomento:} Disequazioni di primo grado}\par\vspace{1mm}


\noindent\textbf{Risposta corretta:} A.\par\vspace{1mm}


\begin{enumerate}[label=\arabic*)]

  \item Isola il termine con l'incognita: $3x = 10 + 8 = 18$.

  \item Dividi entrambi i membri per $3$: $x = \dfrac{18}{3} = 6$.

  \item Verifica: $3\cdot 6 - 8 = 10$, quindi la soluzione è corretta.

\end{enumerate}
\noindent\textbf{Criteri di punteggio:}
\begin{itemize}

  \item 4\{,\}5 pt risposta corretta

  \item 0 pt risposta errata o non data

\end{itemize}

\erroricomuni{Errore di segno nel trasporto di un termine; Divisione non eseguita per il coefficiente dell'incognita; Verifica non effettuata}


\needspace{9\baselineskip}
\subsection*{Quesito 2\hfill\normalsize\mdseries [4,5~pt]}

{\small \textit{Argomento:} Polinomi e scomposizioni}\par\vspace{1mm}


\noindent\textbf{Risposta corretta:} B.\par\vspace{1mm}


\begin{enumerate}[label=\arabic*)]

  \item Cerchiamo due numeri con somma $-2$ e prodotto $1$.

  \item I numeri sono $-1$ e $-1$.

  \item Quindi $x^2 + 2x + 1 = $(x + 1)(x + 1)\$\$.

\end{enumerate}
\noindent\textbf{Criteri di punteggio:}
\begin{itemize}

  \item 4\{,\}5 pt risposta corretta

  \item 0 pt risposta errata o non data

\end{itemize}

\erroricomuni{Segni dei fattori invertiti; Scambiare somma e prodotto; Dimenticare il coefficiente del termine noto}



\titolosezione{Quesiti a risposta breve}

\needspace{9\baselineskip}
\subsection*{Quesito 3\hfill\normalsize\mdseries [6,75~pt]}

{\small \textit{Argomento:} Insiemi numerici e percentuali}\par\vspace{1mm}



\begin{enumerate}[label=\arabic*)]

  \item Sconto: $190 \cdot \dfrac{50}{100} = 95$ euro.

  \item Prezzo finale: $190 - 95 = 95$ euro.

  \item Verifica: il prezzo finale è il 50\% di quello iniziale.

\end{enumerate}
\noindent\textbf{Criteri di punteggio:}
\begin{itemize}

  \item 2\{,\}75 pt impostazione corretta e dati individuati

  \item 2\{,\}75 pt procedimento e calcoli corretti

  \item 1\{,\}25 pt risultato e verifica

\end{itemize}

\erroricomuni{Calcolare lo sconto e fermarsi; Sottrarre la percentuale invece della sua parte; Errore nel riportare gli euro}


\needspace{9\baselineskip}
\subsection*{Quesito 4\hfill\normalsize\mdseries [6,75~pt]}

{\small \textit{Argomento:} Sistemi lineari}\par\vspace{1mm}



\begin{enumerate}[label=\arabic*)]

  \item Isola il termine con l'incognita: $2x = 5 - 1 = 4$.

  \item Dividi entrambi i membri per $2$: $x = \dfrac{4}{2} = 2$.

  \item Verifica: $2\cdot 2 + 1 = 5$, quindi la soluzione è corretta.

\end{enumerate}
\noindent\textbf{Criteri di punteggio:}
\begin{itemize}

  \item 2\{,\}75 pt impostazione corretta e dati individuati

  \item 2\{,\}75 pt procedimento e calcoli corretti

  \item 1\{,\}25 pt risultato e verifica

\end{itemize}

\erroricomuni{Errore di segno nel trasporto di un termine; Divisione non eseguita per il coefficiente dell'incognita; Verifica non effettuata}



\titolosezione{Esercizi}

\needspace{9\baselineskip}
\subsection*{Quesito 5\hfill\normalsize\mdseries [13,75~pt]}

{\small \textit{Argomento:} Funzioni e proporzionalità}\par\vspace{1mm}



\begin{enumerate}[label=\arabic*)]

  \item Sostituisci $x = 2$ nell'espressione: $y = 3\cdot 2 - 4$.

  \item Calcola: $y = 6 - 4 = 2$.

  \item Il punto $\left(2; 2\right)$ appartiene alla retta.

\end{enumerate}
\noindent\textbf{Criteri di punteggio:}
\begin{itemize}

  \item 5\{,\}5 pt impostazione corretta e dati individuati

  \item 5\{,\}5 pt procedimento e calcoli corretti

  \item 2\{,\}75 pt risultato e verifica

\end{itemize}

\erroricomuni{Errore nel segno del coefficiente angolare; Sostituire il valore di y al posto di x}


\needspace{9\baselineskip}
\subsection*{Quesito 6\hfill\normalsize\mdseries [13,75~pt]}

{\small \textit{Argomento:} Probabilità e statistica descrittiva}\par\vspace{1mm}



\begin{enumerate}[label=\arabic*)]

  \item Casi favorevoli: 3.

  \item Casi possibili: $3 + 2 = 5$.

  \item Probabilità: $\dfrac{3}{5}$.

\end{enumerate}
\noindent\textbf{Criteri di punteggio:}
\begin{itemize}

  \item 5\{,\}5 pt impostazione corretta e dati individuati

  \item 5\{,\}5 pt procedimento e calcoli corretti

  \item 2\{,\}75 pt risultato e verifica

\end{itemize}

\erroricomuni{Usare solo le palline di un colore; Non semplificare la frazione quando richiesto}



\titolosezione{Problemi e dimostrazioni}

\needspace{9\baselineskip}
\subsection*{Quesito 7\hfill\normalsize\mdseries [31,75~pt]}

{\small \textit{Argomento:} Equazioni di primo grado}\par\vspace{1mm}



\begin{enumerate}[label=\arabic*)]

  \item Isola il termine con l'incognita: $5x = 16 + 9 = 25$.

  \item Dividi entrambi i membri per $5$: $x = \dfrac{25}{5} = 5$.

  \item Verifica: $5\cdot 5 - 9 = 16$, quindi la soluzione è corretta.

\end{enumerate}
\noindent\textbf{Criteri di punteggio:}
\begin{itemize}

  \item 12\{,\}75 pt impostazione corretta e dati individuati

  \item 12\{,\}75 pt procedimento e calcoli corretti

  \item 6\{,\}25 pt risultato e verifica

\end{itemize}

\erroricomuni{Errore di segno nel trasporto di un termine; Divisione non eseguita per il coefficiente dell'incognita; Verifica non effettuata}


\needspace{9\baselineskip}
\subsection*{Quesito 8\hfill\normalsize\mdseries [18,25~pt]}

{\small \textit{Argomento:} Geometria del piano: perimetro e area}\par\vspace{1mm}



\begin{enumerate}[label=\arabic*)]

  \item Area: $A = b \cdot h = 4 \cdot 9 = 36\,\mathrm{cm}^2$.

  \item Perimetro: $P = 2(b + h) = 2(4 + 9) = 26\,\mathrm{cm}$.

  \item Controllo: il perimetro è maggiore del doppio della base, come atteso.

\end{enumerate}
\noindent\textbf{Criteri di punteggio:}
\begin{itemize}

  \item 7\{,\}25 pt impostazione corretta e dati individuati

  \item 7\{,\}25 pt procedimento e calcoli corretti

  \item 3\{,\}75 pt risultato e verifica

\end{itemize}

\erroricomuni{Confondere area e perimetro; Dimenticare il fattore 2 nel perimetro; Unità di misura errate}




\vspace{6mm}
\noindent\small\textit{Documento generato da VerificheLab. La numerazione e i punteggi corrispondono alla verifica.}

\end{document}
